All numbers in this section come from the \texttt{full} profile (factor-2 size
grids, 5 timed repetitions after 2 warm-ups, 2\,411 cells, 5.3 hours) on the
machine of Section~\ref{sec:method}, supplemented by the \texttt{refine}
profile (factor-1.19 grids around the crossings) and the \texttt{patch}
profile described in Section~\ref{sec:method-patch}.
Table~\ref{tab:crossovers} gives every crossing with its interval and class.

\subsection{Sorting: radix vs.\ introsort}

\begin{figure}[t]
  \centering
  \includegraphics[width=\linewidth]{fig_runtime_sorting}
  \caption{Median runtime of LSD radix sort (11-bit digits) and introsort by
  input family. Dotted lines mark the $n$ at which the working set
  ($16n$ bytes) equals L1d, L2 and L3. Radix sort sits on a flat floor
  below $n\approx10^3$, where the histogram pass dominates.}
  \label{fig:runtime-sorting}
\end{figure}
\begin{figure}[t]
  \centering
  \includegraphics[width=0.72\linewidth]{fig_ratio_sorting_by_dist}
  \caption{Runtime ratio radix/introsort by input family (11-bit digits).
  Uniform keys of four widths cross between 1.4\,k and 1.7\,k; nearly-sorted
  and reverse-sorted input never cross.}
  \label{fig:ratio-sorting}
\end{figure}

On uniformly random keys radix sort wins from about 1\,500 elements and the
crossing is sharp (Figures~\ref{fig:runtime-sorting}
and~\ref{fig:ratio-sorting}): 1.43\,k keys for 8-bit values, 1.60\,k for
16-bit, 1.51\,k for 32-bit and 1.71\,k for 64-bit, each with an interval of
$\pm10\%$ or less. The near-independence from key width is itself a finding:
with 11-bit digits a 64-bit key needs six passes and an 8-bit key one, yet the
crossing barely moves, because below $\sim10^3$ keys both algorithms are
dominated by fixed costs---the histogram and prefix-sum setup on one side,
insertion-sort leaves on the other---and above it both scale as $n$ to within
a few percent over the whole range (fitted exponents 1.00--1.06 for radix,
1.02--1.14 for introsort). Above the crossing the advantage settles between
$2.7\times$ (64-bit keys: $2.44$\,s against $6.63$\,s at $2^{26}$) and
$4\times$ (16-bit keys: $1.14$\,s against $4.55$\,s).

Digit width is the practitioner's knob and it is not free. Eleven-bit digits
are best at every key width; eight-bit digits cost one or two extra passes
and move the 64-bit crossing only to 1.67\,k, but sixteen-bit digits are a
mistake for wide keys: the $2^{16}$-counter histogram per pass is 512\,KB,
larger than L1d, and at 64-bit keys the crossing moves to 7\,k and the
asymptotic advantage shrinks to $1.4\times$ (the gradual interval reaching
8.6\,M in Appendix~\ref{app:table} reflects the ratio hovering between 0.5
and 0.75 across the whole range).

The structured families overturn the verdict. With only 100 distinct values
introsort's partitioning handles duplicates well and the crossing moves to
2.7\,k with an eventual advantage of only $2\times$. On nearly-sorted input
radix sort \emph{never} wins: introsort is $3.7\times$ faster at $2^{26}$ keys
($1.15$\,s vs $4.26$\,s). On reverse-sorted input it is $8\times$ faster
($0.60$\,s vs $4.92$\,s)---median-of-three pivot selection turns a reversed
array into perfectly balanced partitions with perfectly predictable branches,
and introsort runs $11\times$ faster on it than on random data of the same
size. The organ-pipe input, ascending then descending, is the intermediate
case: radix sort crosses at 1.6\,k but only ever gets $1.4\times$ ahead.
Radix sort performs the same scattered writes whether the input is sorted or
not; an algorithm that is oblivious to structure cannot exploit it.

\subsection{Matrix multiplication: Strassen vs.\ naive}

\begin{figure}[t]
  \centering
  \includegraphics[width=0.85\linewidth]{fig_heatmap_matmul}
  \caption{$\log_2$ of the runtime ratio Strassen/naive ($i$-$k$-$j$) over
  matrix size and base-case cutoff on the power-of-two-rich grid; blue means
  Strassen wins.}
  \label{fig:heatmap-matmul}
\end{figure}

Strassen's algorithm is practical much earlier than folklore suggests, with
one important qualification about matrix shape. On sizes with many factors of
two ($n = 2^k$ and $3\cdot2^k$, the grid every benchmark uses) and a base-case
cutoff of 32, it beats the $i$-$k$-$j$ triple loop at every size from $n=32$
up---$0.83\times$ at $n=64$, $0.71\times$ at $256$, $0.39\times$ at $2048$,
$0.31\times$ at $4096$ ($13.4$\,s vs $43.3$\,s)---and the fitted exponent
over the top decade is $2.78$ against the theoretical $2.81$, with the naive
loop at $3.21$.

The refine grid, which includes sizes such as 129, 154, 217 and 259, tells a
different story. Our implementation handles an odd dimension at any level of
the recursion by padding that level to even size, which is the textbook
approach; an $n$ with few factors of two is padded at level after level, and
every padding copies both operands. On those sizes Strassen with cutoff 32 is
up to $1.2\times$ \emph{slower} than naive until $n\approx260$, and the
sustained crossing over all measured sizes is $n\approx300$ (interval
$[231, 308]$), five times later than on the composite grid. Both numbers are
right; which one applies depends on whether the caller controls the matrix
dimensions. Production implementations pad or peel once at the top level,
which would remove this effect, and the comparison shows how much it is worth.

The cutoff behaves as expected (Figure~\ref{fig:heatmap-matmul}): 8 is one
level too deep and costs a factor of $1.6$ at $n=4096$ ($21.2$\,s vs
$13.4$\,s), with the crossing pushed out to between $780$ and $1\,300$ on
uniform matrices (gradual) and to $811$ on integer ones; 16 is $1.2\times$
slower than naive at $n=32$; 32 and 64 are equivalent within noise at the top
of the sweep; 128 and 256 forgo the advantage below 512 and are $30$--$35\%$
slower than cutoff 32 at the top. Integer-valued matrices behave the same as uniform real ones to within
noise, as they should for an algorithm whose work does not depend on the
values.

The second baseline is the cautionary example. The textbook $i$-$j$-$k$ loop
order, which walks the second operand down a column with stride $n$, has a
fitted exponent of $4.3$ over $n\in[229, 2048]$---it is not cubic in practice,
for reasons we take up in Section~\ref{sec:mechanisms}---and at $n=2048$ it
takes $84$\,s where $i$-$k$-$j$ takes $5.0$\,s and Strassen $1.9$\,s. Swapping
two loop headers in the \emph{simple} algorithm is worth $17\times$ at that
size; Strassen's exponent, applied on top, is worth another $2.6\times$.

\subsection{Integer multiplication: Karatsuba vs.\ schoolbook}

\begin{figure}[t]
  \centering
  \includegraphics[width=0.85\linewidth]{fig_heatmap_bigint}
  \caption{$\log_2$ of the runtime ratio Karatsuba/schoolbook over operand
  length (32-bit limbs) and cutoff.}
  \label{fig:heatmap-bigint}
\end{figure}

Karatsuba is the cleanest case. With a cutoff of 32 limbs it overtakes
schoolbook multiplication at 34 limbs (interval $[32, 42]$; roughly $1\,100$
bits), essentially as soon as it does anything at all, and the measured
exponents over the top decade are $1.55$--$1.60$ for Karatsuba against
$1.97$--$2.01$ for schoolbook, within $0.04$ of $\log_2 3$ and $2$. Schoolbook
exceeded the 90-second budget at $4.5\times10^5$ limbs; at the last common
size, $3\times10^5$ limbs ($9.7$\,Mbit operands), Karatsuba is $18\times$
faster ($3.6$\,s vs $65$\,s). The crossing at 34 32-bit limbs is 17 64-bit
limbs (about 1\,090 bits). GMP's tuned parameter file for Alder Lake, the
previous generation of the same core design, sets
\texttt{MUL\_TOOM22\_THRESHOLD} to 13 64-bit limbs (830
bits)~\cite{gmp2022alderlake}, so the two land within a factor of $1.3$ of
each other. The agreement should not be over-read: GMP measures its threshold
against a hand-written assembly basecase and with an assembly Karatsuba, so
both sides of its comparison differ from ours, and the near-coincidence says
mainly that the crossover is insensitive to the quality of the
implementation when both sides are implemented with equal care.

With all-ones limbs---maximal carry propagation---schoolbook is unchanged but
Karatsuba is $1.8\times$ faster than on random limbs, because the sub-products
of the sum terms have predictable carries; the crossing is essentially unchanged (Karatsuba is ahead from the
first size at which it differs from schoolbook). The
cutoff sensitivity is asymmetric (Figure~\ref{fig:heatmap-bigint}): cutoffs of
32, 64 and 128 are indistinguishable at large $n$, 16 makes the crossing
gradual (the two are within noise from 34 to about 60 limbs), and 8 pushes it out
to $\approx700$ limbs and costs a factor of $1.9$ at every size above
($42.8$\,s vs $22.7$\,s at $10^6$ limbs). A cutoff four times too large costs
nothing; four times too small costs $2\times$.

\subsection{Shortest paths: Fibonacci heap vs.\ binary heap}

\begin{figure}[t]
  \centering
  \includegraphics[width=0.72\linewidth]{fig_ratio_dijkstra_by_dist}
  \caption{Runtime ratio Fibonacci-heap/binary-heap (decrease-key) Dijkstra
  by graph family (random graphs: average out-degree 16). Only the complete
  graph reaches the line, and only at the largest size that fits in memory.}
  \label{fig:ratio-dijkstra}
\end{figure}

The Fibonacci heap is the pair that does not cross. On random graphs of
average degree 2, 4, 16, 64 and 256 it trails the indexed binary heap by
$2.1\times$, $1.7\times$, $1.4\times$, $1.3\times$ and $1.3\times$ at the
largest sizes measured ($8.4\times10^6$ vertices at degrees 2--16, $134$\,M
edges at degree 16; $4.2\times10^6$ at degree 64 and $1.0\times10^6$ at 256), on 2-D grids by $2.6\times$ (in line with the $3.15\times$ that
Larkin, Sen and Tarjan measured against an implicit 4-ary heap on the USA road
network~\cite{larkin2014backtobasics}), and the gap narrows only slowly with
$n$ (Figure~\ref{fig:ratio-dijkstra}). The reason is in the operation counts
the harness logs: Dijkstra performs $0.05$ decrease-keys per vertex on
degree-2 graphs, $0.3$ on degree-4 graphs and grids, and $1.4$ at degree 16,
so the $\bigO(1)$ amortised decrease-key that justifies the Fibonacci heap is
exercised far less often than the $\bigO(\log n)$ delete-min that both heaps
share; Section~\ref{sec:mechanisms} argues that the Fibonacci heap's
delete-min is the more expensive of the two.

On complete graphs ($m = n(n-1)$, $6$ decrease-keys per vertex) the ratio
falls steadily from $2.4$ at 100 vertices to $1.0$ at $6\,500$ and $0.98$ at
$7\,700$ vertices ($60$\,M edges, the largest that fits in memory). This is
the regime $m \gg n\log n$ in which $\bigO(m+n\log n)$ should beat
$\bigO(m\log n)$, and the trend is as theory predicts---but what it reaches
within the measurable range is parity, not a win, and the atlas reports it as
such.

Even that parity is against the wrong baseline. The lazy variant,
\texttt{std::priority\_queue} with stale entries skipped on pop, does no
decrease-keys at all and is the fastest of the three on every random graph and
on the complete graph: $78$\,ms against $135$\,ms (Fibonacci) and $138$\,ms
(indexed binary) at $7\,700$ dense vertices; $6.3$\,s against $15.0$\,s and
$8.6$\,s at $8.4$\,M vertices of degree 4---the same ``at least $1.4\times$'' advantage for
dropping decrease-key that Chen et al.\ report on random graphs of average
degree 8~\cite{chen2007priority}. Its heap holds duplicate entries and does
more pops, but each pop touches a contiguous array and the extra work costs
less than any form of decrease-key bookkeeping. Only on grids, where the
frontier is small and duplicates accumulate, does the indexed binary heap edge
it out ($1.89$\,s vs $2.28$\,s). The Fibonacci heap never beats the lazy heap
in any cell, trailing it by $1.7\times$ even on the complete graph.

\subsection{Summary}

\begin{table}[t]
  \centering\scriptsize
  \caption{Sustained crossover points with 95\% bootstrap intervals (shortest
  interval over 1000 resamples of the repetitions). Sharp crossings are given
  as a point; gradual ones (interval wider than a grid step or a factor 1.5)
  as the range over which the two algorithms are within noise of each other;
  ``parity only'' marks a crossing that is never more than 10\% ahead by the
  largest size measured. $\rho(n_{\max})$ is the runtime ratio at the largest
  size measured. For matrices the crossing is reported on the
  power-of-two-rich grid and on all sizes separately. This table shows the
  default parameter of each pair (11-bit digits, cutoff 32); every
  parameter is in Appendix~\ref{app:table}.}
  \label{tab:crossovers}
  \input{../figures/table_crossovers}
\end{table}

Three of the four theoretically faster algorithms win, and early: at
$1\,500$ keys, 32 rows, and 34 limbs---and for the well-tuned configurations
these are sharp crossings with intervals of $\pm10\%$. Each has conditions:
radix sort needs the input not to be nearly or reverse sorted, Strassen needs the matrix
dimension to be rich in factors of two or the implementation to pad only
once, Karatsuba needs its cutoff not to be too small. The fourth reaches
parity on complete graphs at the edge of memory and loses everywhere else,
including to the lazy idiom that practitioners actually write.
